\documentclass[a4paper,12pt]{article}
\usepackage[T2A]{fontenc}
\usepackage[utf8]{inputenc}
\usepackage[russian]{babel}
\usepackage{amsmath,amssymb,mathtools}
\usepackage{helvet}
\renewcommand{\sfdefault}{cmss}
\renewcommand{\familydefault}{\sfdefault}
\usepackage{icomma}
\usepackage[margin=23mm,top=21mm,bottom=23mm,headheight=24pt]{geometry}
\usepackage{microtype}
\usepackage{tikz}
\usetikzlibrary{calc,arrows.meta,positioning,shapes.geometric}
\usepackage{tkz-euclide}
\usepackage{booktabs,tabularx}
\usepackage{enumitem}
\usepackage{fancyhdr}
\usepackage{xcolor}
\usepackage{setspace}
\usepackage[hidelinks]{hyperref}
\definecolor{accent}{HTML}{008080}
\definecolor{darkaccent}{HTML}{00585A}
\definecolor{midgray}{HTML}{555555}
\definecolor{pale}{HTML}{D6E9E7}
\setlength{\parindent}{0pt}
\setlength{\parskip}{5pt}
\setstretch{1.16}
\setlist[enumerate]{leftmargin=*,itemsep=9pt,topsep=6pt,parsep=0pt}
\pagestyle{fancy}
\fancyhf{}
\fancyhead[L]{\footnotesize\color{midgray}Отношения площадей треугольников}
\fancyhead[R]{\footnotesize\color{midgray}Софья Кальней}
\fancyfoot[C]{\footnotesize\thepage}
\renewcommand{\headrulewidth}{0.35pt}
\renewcommand{\headrule}{\hbox to\headwidth{\color{accent}\leaders\hrule height \headrulewidth\hfill}}
\newcommand{\sect}[1]{\vspace{6pt}{\Large\color{darkaccent}#1}\par\vspace{3pt}{\color{accent}\rule{\linewidth}{0.55pt}}\vspace{6pt}}
\newcommand{\subsect}[1]{\vspace{4pt}{\large\color{darkaccent}#1}\par\vspace{2pt}}
\newcommand{\idea}[1]{\textcolor{darkaccent}{\textbf{Главная идея.}} #1}
\newcommand{\good}[1]{\textcolor{darkaccent}{\textbf{Ответ:}} #1}
\begin{document}

% ----- Титульный лист -----
\thispagestyle{empty}
\begin{center}
\vspace*{25mm}
{\large\color{darkaccent} МАТЕМАТИКА}\par
\vspace{23mm}
{\LARGE Чек-лист}\par
\vspace{13mm}
{\LARGE\color{accent} Отношения площадей\\[4pt] треугольников}\par
\vspace{18mm}
{\normalsize Геометрия · Подготовка к ОГЭ}\par
\vfill
{\large Софья Кальней}\par
\vspace{13mm}
{6 ноября 2026 г.}\par
\vspace{22mm}
{\small Лёвин Артём Александрович\\[3pt]эксклюзивно для Софьи Кальней}\par
\vspace{23mm}
\begin{tikzpicture}[x=1cm,y=1cm]
  \draw[accent,line width=0.9pt]
    (0,0) .. controls (1.0,0.50) and (2.3,-0.50) .. (3.4,0)
          .. controls (4.6,0.50) and (5.8,-0.50) .. (6.9,0);
\end{tikzpicture}
\end{center}
\clearpage

\sect{1. Что связывает площади треугольников}
\idea{Если в задаче есть несколько треугольников, найдите у них общий элемент: высоту или угол. Запишите две формулы площади и разделите одну на другую. Общий множитель сократится.}

Площадь треугольника можно выражать двумя способами:
\[
\boxed{S=\frac12 ah}\qquad\text{или}\qquad
\boxed{S=\frac12 ab\sin\gamma},
\]
где $a$~--- сторона, $h$~--- высота к ней; $a,b$~--- две стороны, а $\gamma$~--- угол между ними.

\sect{2. Общая высота: сравниваем основания}
Пусть точка $K$ лежит на стороне $AC$ треугольника $ABC$. У треугольников $ABK$ и $KBC$ высота из вершины $B$ к прямой $AC$ одна и та же.

\begin{center}
\begin{tikzpicture}[scale=0.78]
\tkzDefPoint(0,0){A}
\tkzDefPoint(3.6,3.1){B}
\tkzDefPoint(7.8,0){C}
\tkzDefPoint(5.2,0){K}
\tkzDefPoint(3.6,0){H}
\tkzDrawPolygon[line width=0.85pt](A,B,C)
\tkzDrawSegment[line width=0.95pt,color=accent](B,K)
\tkzDrawSegment[dashed,color=midgray](B,H)
\tkzMarkRightAngle[size=0.25](B,H,C)
\tkzDrawPoints[size=1.6pt](A,B,C,K,H)
\tkzLabelPoints[below](A,H,K,C)
\tkzLabelPoints[above](B)
\end{tikzpicture}
\end{center}

\[
\frac{S_{ABK}}{S_{KBC}}=
\frac{\frac12 AK\cdot BH}{\frac12 KC\cdot BH}
=\boxed{\frac{AK}{KC}}.
\]

\subsect{Пример 1. Известны части основания}
$AK:KC=2:3$, $S_{ABC}=40$. Тогда всё основание соответствует $2+3=5$ частям, поэтому
\[
S_{ABK}=\frac25\cdot40=16,\qquad
S_{KBC}=\frac35\cdot40=24.
\]
\good{$16$ и $24$. Проверка: $16+24=40$.}

\textcolor{darkaccent}{\textbf{Помните:}} если вершина общая, а основания расположены на одной прямой, отношение площадей равно отношению оснований.

\clearpage
\sect{3. Общий угол: сравниваем произведения сторон}
Пусть $M$ лежит на $AB$, а $K$ лежит на $AC$. У треугольников $AMK$ и $ABC$ общий угол $A$.

\begin{center}
\begin{tikzpicture}[scale=0.70]
\tkzDefPoint(0,0){A}
\tkzDefPoint(0,4.6){B}
\tkzDefPoint(7.2,0){C}
\tkzDefPoint(0,2.9){M}
\tkzDefPoint(3.1,0){K}
\tkzDrawPolygon[line width=0.85pt](A,B,C)
\tkzDrawSegment[color=accent,line width=1pt](M,K)
\tkzMarkAngle[size=0.72,color=accent](K,A,M)
\tkzDrawPoints[size=1.6pt](A,B,C,M,K)
\tkzLabelPoints[left](A,B,M)
\tkzLabelPoints[below](K,C)
\end{tikzpicture}
\end{center}

\[
\frac{S_{AMK}}{S_{ABC}}=
\frac{\frac12 AM\cdot AK\sin A}{\frac12 AB\cdot AC\sin A}
=\boxed{\frac{AM}{AB}\cdot\frac{AK}{AC}}.
\]
Важны \emph{доли от целых сторон}. Например, если $AM:MB=3:2$, то $AM/AB=3/5$.

\subsect{Пример 2. Найти площадь внутреннего треугольника}
$AM:MB=3:2$, $AK:KC=2:1$, $S_{ABC}=75$.
Значит, $AM/AB=3/5$ и $AK/AC=2/3$. Получаем
\[
S_{AMK}=\frac35\cdot\frac23\cdot75=30.
\]
\good{$30$.}

\sect{4. Выбираем короткое разбиение фигуры}
Если $A,M,N,B$ расположены на одной прямой именно в таком порядке, а $K$ лежит на $AC$, то
\[
S_{MNK}=S_{AKN}-S_{AKM}.
\]
Оба вычитаемых можно выразить через площадь $ABC$. После вынесения общего множителя:
\[
\boxed{S_{MNK}=S_{ABC}\cdot\frac{MN}{AB}\cdot\frac{AK}{AC}}.
\]

\subsect{Пример 3. Площадь треугольника $MNK$}
$AM:MN:NB=2:2:1$, $AK:KC=1:2$, $S_{ABC}=75$.
Тогда $MN/AB=2/5$, $AK/AC=1/3$ и
\[
S_{MNK}=75\cdot\frac25\cdot\frac13=10.
\]
\good{$10$.}

\clearpage
\sect{5. Четырёхугольник и пересечение диагоналей}
Пусть диагонали $AC$ и $BD$ выпуклого четырёхугольника $ABCD$ пересекаются в точке $O$. Они делят его на четыре треугольника.

\begin{center}
\begin{tikzpicture}[scale=0.71]
\tkzDefPoint(-3.2,0){A}
\tkzDefPoint(0,2.9){B}
\tkzDefPoint(4.5,0){C}
\tkzDefPoint(0,-2.6){D}
\tkzDefPoint(0,0){O}
\tkzDrawPolygon[line width=0.9pt](A,B,C,D)
\tkzDrawSegment[color=accent,line width=0.9pt](A,C)
\tkzDrawSegment[color=accent,line width=0.9pt](B,D)
\tkzDrawPoints[size=1.6pt](A,B,C,D,O)
\tkzLabelPoints[left](A)
\tkzLabelPoints[above](B)
\tkzLabelPoints[right](C)
\tkzLabelPoints[below](D)
\tkzLabelPoint[below right](O){$O$}
\end{tikzpicture}
\end{center}

Треугольники $AOB$ и $BOC$ имеют общую высоту из $B$ к прямой $AC$. Аналогично для треугольников $AOD$ и $DOC$:
\[
\frac{S_{AOB}}{S_{BOC}}=\frac{AO}{OC}
=\frac{S_{AOD}}{S_{DOC}}.
\]
Следовательно, работает важное свойство:
\[
\boxed{S_{AOB}\cdot S_{DOC}=S_{BOC}\cdot S_{AOD}}.
\]
\subsect{Пример 4. Три площади известны}
$S_{AOB}=10$, $S_{BOC}=20$, $S_{AOD}=30$.
Из равенства произведений
\[
10\cdot S_{DOC}=20\cdot30,
\qquad S_{DOC}=60.
\]
Площадь всего четырёхугольника:
\[
S_{ABCD}=10+20+30+60=\boxed{120}.
\]

\sect{6. На что обращать внимание}
\begin{itemize}[leftmargin=*,itemsep=4pt]
\item Если $AM:MB=2:3$, то $AM/AB=2/5$: знаменатель содержит обе части.
\item В формуле через синус берите угол \emph{между выбранными сторонами}.
\item Выбирайте кратчайшее разбиение; проверяйте пары треугольников с общей высотой и разумность ответа.
\end{itemize}

\clearpage
\thispagestyle{fancy}
\begin{center}
{\Large\color{darkaccent} Алгоритм: как сравнивать площади}\par
\vspace{16mm}
\begin{tikzpicture}[
  >=Latex, line width=0.8pt,
  action/.style={draw=darkaccent,rounded corners=2pt,minimum height=1.15cm,text width=5.6cm,align=center,font=\small,inner sep=6pt},
  wide/.style={action,text width=8.0cm},
  decision/.style={draw=darkaccent,diamond,aspect=2.4,text width=3.15cm,align=center,font=\small,inner sep=2pt},
  arrow/.style={-{Latex[length=2.2mm]},color=darkaccent}
]
\node[wide] (start) at (0,0) {Выделите треугольники, площади которых надо связать};
\node[decision] (height) at (0,-2.15) {Есть общая высота?};
\node[action] (hyes) at (-3.8,-4.8) {$\dfrac{S_1}{S_2}=\dfrac{a_1}{a_2}$\\[8pt]сравните основания};
\node[decision] (angle) at (3.85,-4.8) {Есть общий угол?};
\node[action] (ayes) at (3.85,-7.3) {$\dfrac{S_1}{S_2}=\dfrac{a_1b_1}{a_2b_2}$\\[8pt]сравните произведения сторон};
\node[action] (retry) at (3.85,-9.6) {Измените разбиение и повторите поиск общего элемента};
\node[wide] (finish) at (0,-12.15) {Подставьте отношения, вычислите площадь, проверьте ответ};
\draw[arrow] (start) -- (height);
\draw[arrow] (height.west) -| node[pos=0.22,above,font=\small]{да} (hyes.north);
\draw[arrow] (height.east) -| node[pos=0.22,above,font=\small]{нет} (angle.north);
\draw[arrow] (angle.south) -- node[right,font=\small]{да} (ayes.north);
\draw[arrow] (angle.east) -- ++(0.8,0) |- node[pos=0.18,right,font=\small]{нет} (retry.east);
\draw[arrow] (hyes.south) |- (finish.west);
\draw[arrow] (ayes.west) -- (-0.2,-7.3) -- (-0.2,-11.42) -- (finish.north);
\draw[arrow] (retry.east) -- (7.55,-9.6) -- (7.55,0) -- (start.east);
\end{tikzpicture}
\end{center}

\clearpage
\sect{Тренировочный блок · Путь А}
\textcolor{midgray}{10 заданий от простого к сложному. Решайте по чертежу, если он требуется, и записывайте отношения площадей перед вычислениями.}

\begin{enumerate}
\item В треугольнике $ABC$ точка $K$ лежит на стороне $AC$, причём $AK:KC=2:5$. Площадь $ABC$ равна $42$. Найдите площадь $ABK$.

\item В треугольнике $ABC$ точка $D$ лежит на стороне $AB$, причём $AD:DB=3:2$. Площадь $ACD$ равна $21$. Найдите площадь $BCD$.

\item Точка $K$ лежит на стороне $AC$ треугольника $ABC$. Известно, что $S_{ABK}=15$ и $S_{KBC}=25$. Найдите отношение $AK:KC$.

\item На стороне $AB$ треугольника $ABC$ расположены точки $M$ и $N$ в порядке $A,M,N,B$. Известно, что $AM:MN:NB=1:2:3$, а $S_{ABC}=72$. Найдите $S_{CMN}$.

\item В треугольнике $ABC$ точки $M\in AB$, $K\in AC$. Известно, что $AM:MB=1:3$, $AK:KC=2:3$, а $S_{ABC}=100$. Найдите $S_{AMK}$.

\item В треугольнике $ABC$ точки $N\in AB$, $K\in AC$. Известно, что $AN:NB=2:3$, $AK:KC=3:2$, а $S_{ABC}=75$. Найдите $S_{AKN}$.

\item Точки $M,N$ лежат на стороне $AB$ треугольника $ABC$ в порядке $A,M,N,B$; точка $K$ лежит на $AC$. Если $AM:MN:NB=2:2:1$, $AK:KC=1:2$, $S_{ABC}=90$, найдите $S_{MNK}$.

\item В треугольнике $ABC$ точки $M,N$ расположены на $AB$ в порядке $A,M,N,B$, а $K$ лежит на $AC$. Известно, что $AM:MN:NB=1:2:1$, $AK:KC=3:1$, $S_{ABC}=120$. Найдите $S_{MNK}$.

\item Диагонали четырёхугольника $ABCD$ пересекаются в точке $O$. Известно, что $S_{AOB}=12$, $S_{BOC}=18$, $S_{AOD}=20$. Найдите площадь четырёхугольника $ABCD$.

\item Диагонали выпуклого четырёхугольника разделили его на четыре треугольника. Площади трёх из них равны $12$, $18$ и $30$. Площадь четвёртого треугольника больше каждой из указанных трёх. Найдите площадь четырёхугольника.
\end{enumerate}

\clearpage
\sect{Ответы к тренировочному блоку}
\vspace{8pt}
\begin{center}
\renewcommand{\arraystretch}{1.7}
\begin{tabular}{@{}>{\bfseries}c p{8cm}@{}}
\toprule
№ задания & Ответ \\
\midrule
1 & $12$ \\
2 & $14$ \\
3 & $3:5$ \\
4 & $24$ \\
5 & $10$ \\
6 & $18$ \\
7 & $12$ \\
8 & $45$ \\
9 & $80$ \\
10 & $105$ \\
\bottomrule
\end{tabular}
\end{center}
\vspace{12mm}
\sect{Итоговая памятка}
\renewcommand{\arraystretch}{1.55}
\begin{tabularx}{\linewidth}{@{}p{0.31\linewidth}X@{}}
\toprule
\textbf{Что заметили} & \textbf{Что сделать}\\
\midrule
Общая высота & Разделить площади; оставить отношение оснований.\\
Общий угол & Разделить площади; оставить отношение произведений сторон.\\
Точки на одной стороне & Перейти от отношения частей к доле всей стороны.\\
Фигура составлена из частей & Подобрать самое короткое сложение или вычитание площадей.\\
Диагонали четырёхугольника & Использовать $S_{AOB}S_{DOC}=S_{BOC}S_{AOD}$.\\
\bottomrule
\end{tabularx}
\end{document}
